\chapter{Základné pojmy}

Kybernetická a informačná bezpečnosť (KIB) je multidisciplinárna oblasť ľudskej činnosti, zaoberajúca sa skúmaním
hrozieb voči informáciám, informačným systémom a sieťam a návrhom riešení, ktoré by naplneniu hrozieb zamedzili,
alebo aspoň zmiernili ich dopad. KIB je relatívne nová a dynamicky sa rozvíjajúca oblasť, v ktorej sa terminológia ešte
neustálila\footnote{\acrshort{ENISA} publikovala štúdiu, v ktorej porovnávala rozlične výklady pojmu kybernetická
bezpečnosť naprieč šiestimi rôznymi inštitúciami.}, slovenská terminológia sa zväčša prekladá alebo preberá termíny
z anglického jazyka. V tejto kapitole definujeme len základné pojmy a úlohy, ktoré rieši KIB. Budeme vychádzať zo slovníka a
krátkeho úvodu do problematiky \acrshort{KIB}\cite{KUDIKB-MVS}. Ak je čitateľ znalý v problematike \acrshort{KIB} môže
túto časť preskočiť pokojne prejsť k ďalšej kapitole. 


%Budeme vychádzať z úvodnej kapitoly dokumentu \cite{KomentarZoKB}. Ak by mal čitateľ
%záujem o podrobnejšie informácie, dávame do pozornosti krátky úvod do problematiky \acrshort{KIB}\cite{KUDIKB-MVS}.

\begin{itemize}
  \item \textbf{aktívum}, angl. \textit{asset}, je čokoľvek, čo má pre organizáciu hodnotu, môžu byť cieľom útoku a má zmysel 
    (resp. je potrebné) chrániť pred potenciálnymi hrozbami. Aktíva sú hmotné (zariadenia, infraštruktúra, personál) a nehmotné
    (peniaze, informácie, vedomosti).

  \item \textbf{primárne aktíva}, sú aktivity zamarané na plnenie poslania organizácie. 

  \item \textbf{sekundárne aktíva}, sú prostriedky, ktoré umožňujú alebo podporujú činnosť primárnych aktív. 

  \item \textbf{hrozba}, angl. \textit{threat}, je objektívne existujúca potenciálna možnosť priamo alebo nepriamo narušiť
    (alebo negatívne ovplyvniť) systém, spracovávané informácie, alebo iné aktíva organizácie.

  \item \textbf{dôvernosť}, angl. \textit{confidentiality}, je bezpečnostná požiadavka, ktorej naplnenie znamená, že informáciu
    obsiahnutú v správe sa nedozvedia nepovolené osoby.

  \item \textbf{integrita}, angl. \textit{integrity}, je bezpečnostná požiadavka, ktorej naplnenie znamená, že údaje nie je
    možne zmeniť bez toho, aby to ich vlastník alebo adresát nevedel zistiť.

  \item \textbf{dostupnosť}, angl. \textit{availability}, je bezpečnostná požiadavka, ktorej naplnenie znamená, že zdroje systému
    sú k dispozícií oprávnenej - osobe v momente, keď o to požiada; alebo od istého okamihu; alebo v istom časovom rámci
    (napr. dni v týždni, časový interval počas dňa).

  \item \textbf{riziko}, angl. \textit{risk}, je veličina závislá od závažnosti (možného dopadu) hrozby a pravdepodobnosti, 
    že sa hrozba naplní.

  \item \textbf{dopad hrozby}, angl. \textit{threat impact}, negatívne dôsledky naplnenia hrozby na aktívach organizácie.

  \item \textbf{hodnota rizika}, je stredná hodnota dopadu rizika na aktívum, 

  \item \textbf{analýza rizík}, angl. \textit{risk analysis}, proces identifikácie rizík, ohodnotenia rizík a stanovenia
    úrovne rizík.
  \item \textbf{incident}, angl. \textit{incident}, je udalosť alebo situácia, ktorá spôsobí alebo môže spôsobiť nežiadúce
    prerušenie činnosti, stratu, núdzový stav. zranenie, usmrtenie alebo krízu v organizácií alebo systéme.

  \item \textbf{kybernetická a informačná bezpečnosť} je multidisciplinárna oblasť, ktorá skúma hrozby voči aktívam a
    hľadá opatrenia, ktoré majú eliminovať riziká vyplývajúce z hrozieb voči aktívam. Táto disciplína skúma metódy
    dlhodobej, nákladovo a funkčne efektívnej ochrany aktív organizácie.

  \item \textbf{systém riadenia informačnej bezpečnosti}, angl. \textit{\acrfull{ISMS}}, je súbor nástrojov, metód a postupov pre zaistenie
    systematického riešenia informačnej bezpečnosti v organizácií. Cieľom tohto systému je zaručiť dôvernosť, integritu a dostupnosť
    informácie, procesov a systémov informačných technológií. Podrobnejšie sa ním budeme zaoberať v kapitole \ref{chapter:isms}
    o zavádzaní ISMS.

\end{itemize}

Ďalšie pojmy KIB, s ktorými sa stretneme, priebežne vysvetlíme. Čitateľovi, ktorý by mal záujem o terminológiu KIB, odporúčame
preštudovať si slovník\cite{KUDIKB-MVS}.
